


Your profile page is the central view of your student record in UniZ. It shows your personal details, academic information, family details, and your outpass/outing history.

---

\subsection{Profile information}

Your profile is divided into tabs:


\begin{tcolorbox}[colback=gray!5,colframe=primary,title=Documentation Note]

  
\begin{tcolorbox}[colback=gray!5,colframe=primary,title=Documentation Note]

    Full name, gender, blood group, phone number, and date of birth.
  
\end{tcolorbox}

  
\begin{tcolorbox}[colback=gray!5,colframe=primary,title=Documentation Note]

    Student ID (roll number), year, branch, section, and room number.
  
\end{tcolorbox}

  
\begin{tcolorbox}[colback=gray!5,colframe=primary,title=Documentation Note]

    Father's and mother's name, occupation, email, and phone number.
  
\end{tcolorbox}

  
\begin{tcolorbox}[colback=gray!5,colframe=primary,title=Documentation Note]

    Outpass and outing request history, and the option to submit new requests.
  
\end{tcolorbox}

</CardGroup>

**API response example — your profile:**

\begin{lstlisting}[language=json]
{
  "success": true,
  "student": {
    "_id": "550e8400-e29b-41d4-a716-446655440000",
    "username": "O210008",
    "name": "DESU SABER",
    "year": "3",
    "branch": "CSE",
    "profile_url": "https://...",
    "has_pending_requests": false,
    "is_in_campus": true
  }
}

\end{lstlisting}

---

\subsection{Update your details}

You can update the following fields from the Personal and Family tabs:

- Phone number
- Address
- Blood group
- Date of birth
- Gender
- Father's name and occupation
- Mother's name and occupation
- Father's phone number
- Mother's phone number


\section{Instruction Step}

  
\section{Instruction Step}

    On the profile header, tap the pencil icon next to your name, or use the
    **Edit** button (desktop: top-right corner).
  
  
\section{Instruction Step}

    Only changed fields are sent to the server — unchanged fields are left
    as-is.
  
  
\section{Instruction Step}

    On desktop, tap **Save** in the top-right corner. On mobile, tap **Save
    Changes** at the bottom of the editing panel.
  
</Steps>

The update endpoint only accepts the fields you change:

\begin{lstlisting}[language=json]
PUT /profile/student/update
{
  "phone": "9876543210",
  "address": "New Address, Ongole"
}

\end{lstlisting}

<Note>
  Academic fields — your roll number, year, branch, section, and room number —
  are managed by the administration and cannot be edited from the portal.
</Note>

---

\subsection{Profile photo}

You can upload or replace your profile photo directly from the profile page.


\section{Instruction Step}

  
\section{Instruction Step}

    Tap the small camera button at the bottom-right of your profile picture.
  
  
\section{Instruction Step}

    Choose an image from your device. Accepted formats are standard image types
    (JPEG, PNG, etc.).
  
  
\section{Instruction Step}

    The photo is uploaded to cloud storage and your profile URL is updated
    automatically. You'll see a loading spinner while it uploads.
  
</Steps>

<Tip>
  Use a clear, recent photo — it may be used for identification purposes in
  campus records.
</Tip>

---

\subsection{Campus presence status}

The `is_in_campus` field on your profile reflects whether you are currently recorded as present on campus.

| Value   | Meaning                                     |
| ------- | ------------------------------------------- |
| `true`  | You are currently on campus                 |
| `false` | You have checked out and are outside campus |

This status is updated by security at the gate:

- When you leave on an approved outpass or outing, security marks you as **out** (`is_in_campus: false`).
- When you return, security marks you as **in** (`is_in_campus: true`).

You cannot change this status yourself. If there's a discrepancy, contact the security office or your warden.

---

\subsection{Request history}

Your outpass and outing history is available in the **Permissions** tab. It shows a paginated list of all past requests with:

- Request type (outpass or outing)
- Date and duration
- Reason you submitted
- Final status (Pending, Approved, Rejected, or Expired)

**API response example:**

\begin{lstlisting}[language=json]
{
  "success": true,
  "history": [
    {
      "_id": "req-123",
      "type": "outpass",
      "status": "pending",
      "reason": "Home visit",
      "requested_time": "2026-02-10T10:00:00Z"
    }
  ],
  "pagination": { "page": 1, "total": 5 }
}

\end{lstlisting}

Results are paginated at 5 per page. Use the pagination controls at the bottom to navigate older records.

---

\subsection{Grievances and support}

If you have a complaint or need help with something not covered in the portal, use the **Help & Support** section accessible from the main navigation.

Submit a brief description of your issue. It will be routed to the appropriate administrative staff.

<Info>
  For urgent matters — such as a rejected leave request you believe is
  incorrect, or a discrepancy in your academic record — contact your caretaker,
  warden, or academic section directly rather than waiting for a grievance
  response.
</Info>


\section{Deep Technical Analysis}
The underlying system architecture for this module is built upon the \textbf{UniZ Microservices Framework}. 
This design ensures that each functional unit (Auth, Profile, Academics) remains isolated, preventing cascading failures across the university ecosystem.

\subsection{Operational Hardening}
Deployments are managed via K3s (Kubernetes) and containerized using high-performance Alpine Linux Docker images. 
Resource management is critical; for instance, the documentation service itself is provisioned with 2GB of RAM to ensure the responsive rendering of complex documentation graphs and state-management components.

\subsection{Enterprise Security Topology}
Every request entering the UniZ gateway is subjected to an aggressive JWT validation layer. 
Permissions are granular, mapping directly onto the RBAC (Role-Based Access Control) matrix defined in the core system specifications.

\subsection{Data Governance}
Prisma-driven data access ensures type safety and predictable query performance. 
We utilize PostgreSQL connection pooling to handle concurrent requests from the student portal and administrative dashboards simultaneously.
